\documentclass[11pt]{article}
% ======================================================================
%  examstyle.tex — EPFL-style exam layout for the weekly Time Series exams
%  Shared by every Exam_Week_NN(.tex / _Solutions.tex). Compiles with pdflatex.
% ======================================================================
\usepackage[utf8]{inputenc}
\usepackage[T1]{fontenc}
\usepackage[english]{babel}
\usepackage[a4paper,margin=2.4cm]{geometry}
\usepackage{amsmath,amssymb,amsthm,mathtools}
\usepackage{bm}
\usepackage{enumitem}
\usepackage{array}

\setlength{\parindent}{0pt}
\setlength{\parskip}{0.5em}
\emergencystretch=3em
\allowdisplaybreaks[1]

% ---------- math macros (same names as the rest of the project) ----------
\newcommand{\E}{\mathbb{E}}
\DeclareMathOperator{\Var}{Var}
\DeclareMathOperator{\Cov}{Cov}
\DeclareMathOperator{\Corr}{Corr}
\DeclareMathOperator*{\argmin}{arg\,min}
\DeclareMathOperator{\MSE}{MSE}
\DeclareMathOperator{\trace}{tr}
\newcommand{\Reals}{\mathbb{R}}
\newcommand{\Ints}{\mathbb{Z}}
\newcommand{\Comp}{\mathbb{C}}
\newcommand{\Nat}{\mathbb{N}}
\newcommand{\Prob}{\mathbb{P}}
\newcommand{\Tr}{^{\mathsf{T}}}
\newcommand{\conj}[1]{\overline{#1}}
\newcommand{\abs}[1]{\left\lvert #1 \right\rvert}
\newcommand{\norm}[1]{\left\lVert #1 \right\rVert}
\newcommand{\ind}{\mathbf{1}}
\newcommand{\eps}{\varepsilon}
\newcommand{\diff}{\nabla}
\newcommand{\acvs}{\gamma}
\newcommand{\acf}{\rho}
\newcommand{\pacf}{\alpha}
\newcommand{\sdf}{\mathcal{S}}
\newcommand{\Bsh}{B}
\newcommand{\iid}{\overset{\text{iid}}{\sim}}

\setlist[enumerate]{leftmargin=2em,topsep=2pt,itemsep=4pt}
\setlist[itemize]{leftmargin=1.6em,topsep=2pt,itemsep=2pt}

\newcounter{exo}

% ---------- exam cover header :  \examheader{exam title}{duration} ----------
\newcommand{\examheader}[2]{%
  \thispagestyle{empty}
  \begin{center}
    {\Large\bfseries Time Series}\\[3pt]
    Spring Semester 2026, EPFL\\
    \'Ecole Polytechnique F\'ed\'erale de Lausanne\\
    Prof.\ Dr.\ Sofia C.\ Olhede\\[7pt]
    \hrule height 0.8pt
    \vspace{9pt}
    {\large\bfseries Practice Exam --- #1}\\[4pt]
    {\normalsize Duration: #2 \quad$\cdot$\quad No calculator \quad$\cdot$\quad Answer on paper}\\
    \vspace{8pt}
    \hrule height 0.8pt
  \end{center}
  \vspace{14pt}
  \begin{tabular}{@{}p{8.2cm}@{}p{8cm}@{}}
    Name:\enspace\hrulefill & Surname:\enspace\hrulefill
  \end{tabular}\\[14pt]
  SCIPER (Student Card Number):\enspace\hrulefill
  \vspace{16pt}
}

% ---------- points table (fixed: 4 problems) ----------
\newcommand{\pointstable}{%
  \begin{center}
  \renewcommand{\arraystretch}{1.5}
  \begin{tabular}{|p{5cm}|p{4cm}|}
    \hline
    \textbf{Exercise} & \textbf{Points}\\\hline
    1 & \\\hline
    2 & \\\hline
    3 & \\\hline
    4 & \\\hline
    \textbf{Total} & \\\hline
  \end{tabular}
  \end{center}
}

% ---------- problem heading (exam: one problem per page) ----------
\newcommand{\problem}[1]{%
  \clearpage
  \stepcounter{exo}%
  \begin{center}\textbf{\large Problem \theexo\ (#1 points)}\end{center}
  \vspace{4pt}
}
% ---------- answer space: fill the statement page + one blank answer page ----------
% Put \answerpage at the END of each problem so every exercise gets its own page(s).
\newcommand{\answerpage}{\par\vfill\newpage\null\newpage}

% ---------- solutions document ----------
\newcommand{\solheader}[1]{%
  \begin{center}
    {\Large\bfseries Time Series --- Solutions}\\[3pt]
    {\large Practice Exam --- #1}\\[2pt]
    EPFL, MATH-342 \quad$\cdot$\quad Prof.\ S.\ C.\ Olhede\\[5pt]
    \hrule height 0.8pt
  \end{center}
  \vspace{10pt}
}
\newcommand{\solproblem}[1]{%
  \bigskip
  \stepcounter{exo}%
  {\large\bfseries Problem \theexo\ (#1 points)}\par\nobreak\vspace{2pt}
}
% Rappel de l'énoncé avant la solution (dans les corrigés)
\newcommand{\statementstart}{\par\vspace{1pt}{\bfseries\itshape Statement.}\par\nobreak\begingroup\small\itshape}
\newcommand{\statementend}{\par\endgroup\vspace{5pt}\hrule\vspace{6pt}{\bfseries Solution.}\par\nobreak\vspace{2pt}}

% --- local macros not in examstyle ---
\newcommand{\Herm}{^{\mathsf{H}}}
\DeclareMathOperator{\diag}{diag}
\newcommand{\Fej}{\mathcal{F}}
\begin{document}

\solheader{Week 9: Multivariate Time Series \& VAR}

% ---------------------------------------------------------------
\solproblem{5}
\statementstart
We work with bivariate series $X_t=(X_t^{(1)},X_t^{(2)})\Tr$. Recall the matrix
autocovariance $\Gamma_\tau=\Cov(X_{t+\tau},X_t)$ (\emph{the lag $\tau$ sits in
the first argument}), with $(j,k)$ entry the cross-covariance
$\acvs_\tau^{(j,k)}=\Cov(X_{t+\tau}^{(j)},X_t^{(k)})$, and the
cross-correlation $\acf_\tau^{(j,k)}=\acvs_\tau^{(j,k)}/\sqrt{\acvs_0^{(j,j)}\acvs_0^{(k,k)}}$.

\textbf{(a)} State the symmetry relation linking $\Gamma_{-\tau}$ to $\Gamma_\tau$, prove it from the definition, and explain in one sentence why the off-diagonal cross-covariance $\acvs_\tau^{(1,2)}$ need \emph{not} be an even function of $\tau$ (unlike a scalar ACVS). \hfill(1.5 pts)

\textbf{(b)} Let $\{U_t^{(1)}\}$ and $\{U_t^{(2)}\}$ be two jointly stationary, mutually independent zero-mean series with $\Var(U_t^{(j)})=\sigma_j^2$, and let $\{\phi_t\}$ be a scalar mean-zero white noise of variance $\sigma_\phi^2$, independent of both. Define the common-shock model
\[
Y_t^{(1)}=U_t^{(1)}+\phi_t,\qquad Y_t^{(2)}=U_t^{(2)}+\phi_t .
\]
Compute the cross-covariance $\acvs_\tau^{(1,2)}$ and the cross-correlation $\acf_\tau^{(1,2)}$ between $\{Y^{(1)}\}$ and $\{Y^{(2)}\}$ at \emph{all} lags $\tau$, and identify at which lag(s) the shared shock injects correlation. \hfill(2 pts)

\textbf{(c)} Specialise to $\sigma_1^2=3$, $\sigma_2^2=1$, $\sigma_\phi^2=1$ and report the numerical value of $\acf_0^{(1,2)}$ and of $\acf_\tau^{(1,2)}$ for $\tau\neq0$. \hfill(1.5 pts)
\statementend

\textbf{(a)} The symmetry relation is
\[
\Gamma_{-\tau}=\Gamma_\tau\Tr
\qquad\Longleftrightarrow\qquad
\acvs_{-\tau}^{(j,k)}=\acvs_\tau^{(k,j)} .
\]
\emph{Proof.} Covariance is symmetric in its two arguments, and the process is stationary, so for every $j,k$,
\[
\acvs_\tau^{(j,k)}=\Cov\bigl(X_{t+\tau}^{(j)},X_t^{(k)}\bigr)
=\Cov\bigl(X_t^{(k)},X_{t+\tau}^{(j)}\bigr)
\overset{s=t+\tau}{=}\Cov\bigl(X_{s-\tau}^{(k)},X_s^{(j)}\bigr)
=\acvs_{-\tau}^{(k,j)} .
\]
Reading this off entrywise: the $(j,k)$ entry of $\Gamma_\tau$ equals the $(k,j)$ entry of $\Gamma_{-\tau}$, i.e.\ $\Gamma_{-\tau}=\Gamma_\tau\Tr$.

\emph{Why not even.} Flipping the sign of $\tau$ forces a \emph{swap of the channel indices} $(j,k)\to(k,j)$; only the diagonal entries ($j=k$) are therefore even in $\tau$. For $j\neq k$ one has merely $\acvs_{-\tau}^{(1,2)}=\acvs_\tau^{(2,1)}$, which differs from $\acvs_\tau^{(1,2)}$ in general --- this asymmetry is exactly what encodes lead--lag (directional) relationships between the two channels.

\textbf{(b)} Expand the cross-covariance by bilinearity, using independence of $U^{(1)},U^{(2)},\phi$ (all cross terms involving distinct independent zero-mean series vanish):
\[
\begin{aligned}
\acvs_\tau^{(1,2)}
&=\Cov\bigl(U_{t+\tau}^{(1)}+\phi_{t+\tau},\ U_t^{(2)}+\phi_t\bigr)\\
&=\underbrace{\Cov\bigl(U_{t+\tau}^{(1)},U_t^{(2)}\bigr)}_{0}
+\underbrace{\Cov\bigl(U_{t+\tau}^{(1)},\phi_t\bigr)}_{0}
+\underbrace{\Cov\bigl(\phi_{t+\tau},U_t^{(2)}\bigr)}_{0}
+\Cov\bigl(\phi_{t+\tau},\phi_t\bigr).
\end{aligned}
\]
The first term is $0$ because $U^{(1)},U^{(2)}$ are independent, the two middle terms because $\phi$ is independent of the $U$'s, and $\Cov(\phi_{t+\tau},\phi_t)=\sigma_\phi^2\,\delta_{\tau,0}$ because $\phi$ is white noise. Hence
\[
\boxed{\;\acvs_\tau^{(1,2)}=\sigma_\phi^2\,\delta_{\tau,0}
=\begin{cases}\sigma_\phi^2 & \tau=0,\\[2pt] 0 & \tau\neq0.\end{cases}\;}
\]
The marginal variances are $\acvs_0^{(1,1)}=\Var(U_t^{(1)}+\phi_t)=\sigma_1^2+\sigma_\phi^2$ and likewise $\acvs_0^{(2,2)}=\sigma_2^2+\sigma_\phi^2$. Therefore
\[
\boxed{\;
\acf_\tau^{(1,2)}
=\frac{\acvs_\tau^{(1,2)}}{\sqrt{(\sigma_1^2+\sigma_\phi^2)(\sigma_2^2+\sigma_\phi^2)}}
=\begin{cases}
\dfrac{\sigma_\phi^2}{\sqrt{(\sigma_1^2+\sigma_\phi^2)(\sigma_2^2+\sigma_\phi^2)}} & \tau=0,\\[12pt]
0 & \tau\neq0.
\end{cases}\;}
\]
The shared shock injects correlation \emph{only at lag $0$} (contemporaneously); at all nonzero lags the two channels are uncorrelated. (This is the ``common shock $\Rightarrow$ everything moves together'' mechanism: $\Gamma_0$ gains the rank-one term $\sigma_\phi^2\,\mathbf 1\mathbf 1\Tr$.)

\textbf{(c)} With $\sigma_1^2=3,\ \sigma_2^2=1,\ \sigma_\phi^2=1$: $\sigma_1^2+\sigma_\phi^2=4$ and $\sigma_2^2+\sigma_\phi^2=2$, so
\[
\acf_0^{(1,2)}=\frac{1}{\sqrt{4\cdot2}}=\frac{1}{\sqrt8}=\frac{1}{2\sqrt2}=\frac{\sqrt2}{4}\approx0.3536,
\qquad
\acf_\tau^{(1,2)}=0\quad(\tau\neq0).
\]

% ---------------------------------------------------------------
\solproblem{5}
\statementstart
\textbf{(The delay process and its cross-spectrum.)} Fix an integer delay $\nu$ and a gain $\alpha>0$. Let $\{X_t^{(2)}\}$ be a zero-mean stationary series with autocovariance $\acvs_\tau^{(2,2)}$, spectral density $\sdf_2(f)=\sum_\tau\acvs_\tau^{(2,2)}e^{-2\pi if\tau}$ and variance $\sigma_2^2=\acvs_0^{(2,2)}$. Define a second channel as a scaled, delayed copy corrupted by noise,
\[
X_t^{(1)}=\alpha\,X_{t-\nu}^{(2)}+\eps_t ,\qquad t\in\Ints,
\]
where $\{\eps_t\}$ (variance $\sigma^2$) is white noise uncorrelated with $\{X_t^{(2)}\}$ at all lags.

\textbf{(a)} Show that the cross-covariance is a \emph{shift} of the autocovariance of $X^{(2)}$, namely $\acvs_\tau^{(1,2)}=\alpha\,\acvs_{\tau-\nu}^{(2,2)}$, and deduce the cross-correlation $\acf_\tau^{(1,2)}$ at all lags in terms of $\acf^{(2,2)}$. \hfill(1.5 pts)

\textbf{(b)} Compute the cross-spectrum $S_{1,2}(f)=\sum_\tau\acvs_\tau^{(1,2)}e^{-2\pi if\tau}$ in closed form, and read off its \textbf{amplitude} $\abs{S_{1,2}(f)}$ and its \textbf{phase} $\arg S_{1,2}(f)$. Why is the phase \emph{linear} in $f$, and what is its slope? \hfill(2 pts)

\textbf{(c)} Define the coherence $r_{1,2}(f)=S_{1,2}(f)/\sqrt{S_{1,1}(f)\,S_{2,2}(f)}$. Using $S_{1,1}(f)=\alpha^2\sdf_2(f)+\sigma^2$ and $S_{2,2}(f)=\sdf_2(f)$, give $\abs{r_{1,2}(f)}$, and show that in the \emph{noiseless} limit $\sigma^2\to0$ the magnitude-squared coherence equals $1$ at every frequency. \hfill(1.5 pts)
\statementend

\textbf{(a)} Since $\{\eps_t\}$ is uncorrelated with $\{X_t^{(2)}\}$ at all lags, $\Cov(\eps_{t+\tau},X_t^{(2)})=0$, so by bilinearity
\[
\acvs_\tau^{(1,2)}=\Cov\bigl(X_{t+\tau}^{(1)},X_t^{(2)}\bigr)
=\alpha\,\Cov\bigl(X_{t+\tau-\nu}^{(2)},X_t^{(2)}\bigr)
+\underbrace{\Cov\bigl(\eps_{t+\tau},X_t^{(2)}\bigr)}_{0}
=\alpha\,\acvs_{\tau-\nu}^{(2,2)} .
\]
The marginal variances are
\[
\acvs_0^{(1,1)}=\Var\bigl(\alpha X_{t-\nu}^{(2)}+\eps_t\bigr)=\alpha^2\sigma_2^2+\sigma^2,
\qquad
\acvs_0^{(2,2)}=\sigma_2^2 ,
\]
(the cross term vanishes by uncorrelatedness). Hence
\[
\acf_\tau^{(1,2)}
=\frac{\alpha\,\acvs_{\tau-\nu}^{(2,2)}}{\sqrt{(\alpha^2\sigma_2^2+\sigma^2)}\,\sqrt{\sigma_2^2}}
=\frac{\acf_{\tau-\nu}^{(2,2)}}{\sqrt{1+\sigma^2/(\alpha\sigma_2)^2}} ,
\]
using $\acvs_{\tau-\nu}^{(2,2)}=\sigma_2^2\,\acf_{\tau-\nu}^{(2,2)}$. So the cross-correlation is the \emph{autocorrelation of $X^{(2)}$ shifted by $\nu$} and shrunk by the constant factor $\bigl(1+\sigma^2/(\alpha\sigma_2)^2\bigr)^{-1/2}\le1$ coming from the added noise.

\textbf{(b)} Fourier-transform the cross-covariance and re-index $\tau\to\tau+\nu$ to pull out the delay phase:
\[
\begin{aligned}
S_{1,2}(f)
&=\sum_{\tau\in\Ints}\acvs_\tau^{(1,2)}e^{-2\pi if\tau}
=\alpha\sum_{\tau\in\Ints}\acvs_{\tau-\nu}^{(2,2)}e^{-2\pi if\tau}\\
&=\alpha\,e^{-2\pi if\nu}\sum_{\tau'\in\Ints}\acvs_{\tau'}^{(2,2)}e^{-2\pi if\tau'}
=\alpha\,\sdf_2(f)\,e^{-2\pi if\nu} .
\end{aligned}
\]
Because $\sdf_2(f)\ge0$ is real (it is the spectral density of $X^{(2)}$) and $\alpha>0$,
\[
\boxed{\;\abs{S_{1,2}(f)}=\alpha\,\sdf_2(f),\qquad
\arg S_{1,2}(f)=-2\pi\nu f .\;}
\]
The phase is \emph{linear in $f$} precisely because a pure time shift of $\nu$ samples in the time domain becomes a multiplicative factor $e^{-2\pi if\nu}$ in the frequency domain (the Fourier shift rule); its slope is $-2\pi\nu$, the textbook signature of a lead/lag of $\nu$ samples. The amplitude is just the spectrum of $X^{(2)}$ scaled by the gain $\alpha$.

\textbf{(c)} With $S_{1,1}(f)=\alpha^2\sdf_2(f)+\sigma^2$, $S_{2,2}(f)=\sdf_2(f)$ and $\abs{S_{1,2}(f)}=\alpha\sdf_2(f)$,
\[
\abs{r_{1,2}(f)}
=\frac{\abs{S_{1,2}(f)}}{\sqrt{S_{1,1}(f)\,S_{2,2}(f)}}
=\frac{\alpha\,\sdf_2(f)}{\sqrt{\bigl(\alpha^2\sdf_2(f)+\sigma^2\bigr)\,\sdf_2(f)}}
=\frac{\alpha\,\sqrt{\sdf_2(f)}}{\sqrt{\alpha^2\sdf_2(f)+\sigma^2}} .
\]
Squaring,
\[
\abs{r_{1,2}(f)}^2=\frac{\alpha^2\sdf_2(f)}{\alpha^2\sdf_2(f)+\sigma^2}\in[0,1].
\]
As $\sigma^2\to0$ this $\to\dfrac{\alpha^2\sdf_2(f)}{\alpha^2\sdf_2(f)}=1$ at every frequency where $\sdf_2(f)>0$: a \emph{pure} delay produces perfect coherence (magnitude $1$) with linear phase $-2\pi\nu f$. The added noise $\eps_t$ is the only thing that pulls the coherence below $1$, and it does so by the factor $\bigl(1+\sigma^2/(\alpha^2\sdf_2(f))\bigr)^{-1}$.

% ---------------------------------------------------------------
\solproblem{5}
\statementstart
\textbf{(VAR stability and the VAR(1) second-order structure.)}

\textbf{(a)} Consider the bivariate $\mathrm{VAR}(2)$ process $X_t=\Phi_1X_{t-1}+\Phi_2X_{t-2}+\eps_t$ with
\[
\Phi_1=\begin{pmatrix}0.5 & 0\\[2pt] 0.3 & 0.7\end{pmatrix},\qquad
\Phi_2=\begin{pmatrix}0 & 0\\[2pt] 0.1 & -0.1\end{pmatrix}.
\]
Form the reverse characteristic polynomial $\det\bigl(I_2-\Phi_1 z-\Phi_2 z^2\bigr)$, exploit the lower-triangular structure to \emph{factor} it, find all its roots, and decide whether the VAR(2) is stable. State the companion-matrix ($F$) form of the stability test and verify it gives the same conclusion. \hfill(2.5 pts)

\textbf{(b)} Now take the \emph{diagonal} $\mathrm{VAR}(1)$ process $X_t=\Phi\,X_{t-1}+\eps_t$ with
\[
\Phi=\diag\!\bigl(\tfrac12,\tfrac13\bigr),\qquad
\{\eps_t\}\sim\mathrm{WN}(0,I_2).
\]
Write its $\mathrm{MA}(\infty)$ representation, and compute the stationary covariance $\Gamma_0$ and the lag-$1$ matrix $\Gamma_1$ in closed form (exact fractions). \hfill(1.5 pts)

\textbf{(c)} For the process of part (b), what is the coherence $r_{1,2}(f)$ between the two channels at every frequency, and why? \hfill(1 pt)
\statementend

\textbf{(a)} Build $M(z)=I_2-\Phi_1 z-\Phi_2 z^2$:
\[
M(z)=\begin{pmatrix}
1-0.5z & 0\\[2pt]
-0.3z-0.1z^2 & 1-0.7z+0.1z^2
\end{pmatrix}.
\]
This matrix is lower triangular, so its determinant is the product of the diagonal entries (the off-diagonal cross-terms $\Phi_1^{(2,1)},\Phi_2^{(2,1)}$ do \emph{not} enter the determinant):
\[
\det M(z)=(1-0.5z)\bigl(1-0.7z+0.1z^2\bigr).
\]
\emph{Roots.} The first factor gives $z=2$. The quadratic $0.1z^2-0.7z+1=0$ has
\[
z=\frac{0.7\pm\sqrt{0.49-0.4}}{0.2}=\frac{0.7\pm0.3}{0.2}
\ \Longrightarrow\ z=2\ \text{ or }\ z=5 .
\]
So the three roots are $z\in\{2,2,5\}$, all with $\abs z>1$ (strictly outside the unit circle). Hence the $\mathrm{VAR}(2)$ is \textbf{stable} (and therefore stationary). \emph{Sanity check:} $\det M(0)=\det I_2=1$, as required.

\emph{Companion form.} Stacking $Y_t=(X_t\Tr,X_{t-1}\Tr)\Tr\in\Reals^4$ turns the VAR(2) into the VAR(1) $Y_t=FY_{t-1}+U_t$ with
\[
F=\begin{pmatrix}\Phi_1 & \Phi_2\\ I_2 & 0\end{pmatrix}
=\begin{pmatrix}
0.5 & 0 & 0 & 0\\
0.3 & 0.7 & 0.1 & -0.1\\
1 & 0 & 0 & 0\\
0 & 1 & 0 & 0
\end{pmatrix}.
\]
The companion test is ``all eigenvalues of $F$ inside the unit circle'', and it is the reciprocal of the root test: $\det(I_2-\Phi_1 z-\Phi_2 z^2)=z^{4}\det(z^{-1}I_4-F)$ up to a constant, so the eigenvalues of $F$ are the reciprocals of the roots $z_i$. Here $\lambda=1/2,\,1/2,\,1/5$ (and a structural $\lambda=0$ from the rank-deficient $\Phi_2$), all with $\abs\lambda<1$, giving the \emph{same} conclusion: stable.

\textbf{(b)} Since $\Phi=\diag(\tfrac12,\tfrac13)$ has spectral radius $\tfrac12<1$, the VAR(1) is stable and back-substitution gives the matrix geometric ($\mathrm{MA}(\infty)$) filter
\[
X_t=\sum_{j=0}^{\infty}\Phi^{\,j}\eps_{t-j},
\qquad
\Phi^{\,j}=\diag\!\bigl(2^{-j},\,3^{-j}\bigr).
\]
By the VAR(1) covariance formula $\Gamma_0=\sum_{k\ge0}\Phi^k\Sigma(\Phi^k)\Tr$ with $\Sigma=I_2$ (everything diagonal, so the two channels decouple into independent AR(1)'s),
\[
\Gamma_0=\sum_{k=0}^{\infty}\diag\!\bigl(4^{-k},\,9^{-k}\bigr)
=\diag\!\left(\frac{1}{1-\tfrac14},\ \frac{1}{1-\tfrac19}\right)
=\boxed{\ \diag\!\left(\tfrac43,\ \tfrac98\right)\ }.
\]
(Equivalently, solve the Lyapunov equation $\Gamma_0=\Phi\Gamma_0\Phi\Tr+I_2$ componentwise: $g_{jj}=\phi_j^2 g_{jj}+1\Rightarrow g_{jj}=1/(1-\phi_j^2)$.) The lag-$1$ matrix follows from Yule--Walker $\Gamma_1=\Phi\Gamma_0$:
\[
\Gamma_1=\diag\!\bigl(\tfrac12,\tfrac13\bigr)\diag\!\bigl(\tfrac43,\tfrac98\bigr)
=\boxed{\ \diag\!\left(\tfrac23,\ \tfrac38\right)\ }.
\]

\textbf{(c)} Both $\Phi$ and $\Sigma=I_2$ are diagonal, so by the formula in (b) \emph{every} $\Gamma_\tau$ is diagonal: all off-diagonal cross-covariances vanish at every lag. Hence the cross-spectrum $S_{1,2}(f)=\sum_\tau\acvs_\tau^{(1,2)}e^{-2\pi if\tau}\equiv0$, and the coherence
\[
r_{1,2}(f)=\frac{S_{1,2}(f)}{\sqrt{S_{1,1}(f)S_{2,2}(f)}}=0\quad\text{for all }f .
\]
The two channels are completely uncoupled: the off-diagonal entries of $\Phi$ (here zero) are precisely the cross-channel dynamics, and there is no shared innovation ($\Sigma$ diagonal), so nothing links the series at any frequency.

% ---------------------------------------------------------------
\solproblem{5}
\statementstart
\textbf{(Revision of Week 8: the tapered periodogram --- bias kernel \& unbiasedness.)} Data $X_1,\dots,X_n$ ($T=\{1,\dots,n\}$, $\Delta=1$) are observed from a zero-mean stationary process with spectral density $\sdf(f)$. A \emph{data taper} is a sequence $\{h_t\}_{t\in T}$ (extended by $h_t=0$ off $T$); the tapered periodogram is
\[
\hat\sdf^{(p)}_h(f)=\abs{J_h(f)}^2 ,\qquad
J_h(f)=\sum_{t\in T}h_t\,X_t\,e^{-2\pi itf},
\]
and $H(f)=\sum_{t}h_t e^{-2\pi itf}$ denotes the discrete Fourier transform of the taper.

\textbf{(a)} Show that the expectation of the tapered periodogram is the frequency convolution
\[
\E\bigl[\hat\sdf^{(p)}_h(f)\bigr]=\int_{-1/2}^{1/2}\sdf(f')\,\abs{H(f-f')}^2\,df' .
\]
(\emph{Hint:} expand $\abs{J_h}^2$ as a double sum, take expectations to get $\sum_\tau w_\tau\acvs_\tau e^{-2\pi if\tau}$ with $w_\tau=\sum_t h_t h_{t+\tau}$, then identify $w$ as an autocorrelation of the taper and apply the convolution theorem.) \hfill(2 pts)

\textbf{(b)} Suppose $\{X_t\}$ is white noise of variance $\sigma^2$ (so $\sdf\equiv\sigma^2$). Using part (a) and Parseval's relation $\int_{-1/2}^{1/2}\abs{H(f')}^2\,df'=\sum_t h_t^2$, prove that $\hat\sdf^{(p)}_h(f)$ is \emph{unbiased} for $\sdf(f)$ at every $f$ and every $n$ \textbf{iff} $\norm{h}_2^2=\sum_t h_t^2=1$. \hfill(1.5 pts)

\textbf{(c)} Take $n=3$ and the un-normalised triangular weights $g=(g_1,g_2,g_3)=(1,2,1)$. Find the constant $c$ so that $h_t=c\,g_t$ is an admissible taper (i.e.\ $\norm h_2^2=1$), and contrast, in one or two sentences, what tapering versus \emph{multitapering} each fix about the raw periodogram. \hfill(1.5 pts)
\statementend

\textbf{(a)} Write $Y_t=X_t$ (zero mean). Expand the squared modulus as a double sum over $T$:
\[
\hat\sdf^{(p)}_h(f)=\abs{J_h(f)}^2
=\Bigl(\sum_{t\in T}h_tX_t e^{-2\pi itf}\Bigr)\Bigl(\sum_{s\in T}h_sX_s e^{2\pi isf}\Bigr)
=\sum_{t\in T}\sum_{s\in T}h_th_s\,X_tX_s\,e^{-2\pi if(t-s)} .
\]
Take expectations, using $\E[X_tX_s]=\acvs_{t-s}$ and grouping the terms by the lag $\tau=t-s$ (with $h$ set to $0$ outside $T$, the sum over $t$ runs over all of $\Ints$):
\[
\E\bigl[\hat\sdf^{(p)}_h(f)\bigr]
=\sum_{t,s}h_th_s\,\acvs_{t-s}\,e^{-2\pi if(t-s)}
=\sum_{\tau\in\Ints} w_\tau\,\acvs_\tau\,e^{-2\pi if\tau},
\qquad
w_\tau=\sum_{t\in\Ints}h_t h_{t+\tau} .
\]
The weight $w_\tau$ is the autocorrelation of the taper. Writing $\tilde h_t=h_{-t}$, we have $w=\tilde h * h$, whose discrete Fourier transform is, by the convolution theorem,
\[
W(f)=\tilde H(f)\,H(f)=\conj{H(f)}\,H(f)=\abs{H(f)}^2
\]
(using $\tilde H(f)=\conj{H(f)}$ for real $h$). Since $\E[\hat\sdf^{(p)}_h(f)]=\sum_\tau w_\tau\acvs_\tau e^{-2\pi if\tau}$ is the transform of the \emph{product} $w_\tau\acvs_\tau$, the convolution theorem (product in time $=$ convolution in frequency) gives
\[
\E\bigl[\hat\sdf^{(p)}_h(f)\bigr]
=\int_{-1/2}^{1/2}\sdf(f')\,W(f-f')\,df'
=\int_{-1/2}^{1/2}\sdf(f')\,\abs{H(f-f')}^2\,df' ,
\]
as claimed. The kernel that blurs the true spectrum is $\abs{H(f)}^2$.

\textbf{(b)} For white noise $\sdf\equiv\sigma^2$, part (a) and periodicity give
\[
\E\bigl[\hat\sdf^{(p)}_h(f)\bigr]
=\int_{-1/2}^{1/2}\sigma^2\,\abs{H(f-f')}^2\,df'
=\sigma^2\int_{-1/2}^{1/2}\abs{H(f')}^2\,df'
=\sigma^2\sum_{t\in T}h_t^2=\sigma^2\,\norm h_2^2 ,
\]
where the third equality is Parseval's relation. This holds for \emph{every} $f$ and every $n$. Therefore
\[
\E\bigl[\hat\sdf^{(p)}_h(f)\bigr]=\sigma^2=\sdf(f)
\quad\Longleftrightarrow\quad
\norm h_2^2=\sum_{t\in T}h_t^2=1 .
\]
So the normalisation $\norm h_2^2=1$ is exactly the condition that makes the tapered periodogram unbiased for a flat (white-noise) spectrum --- this pins down the normalisation in the definition of a taper. (For the constant taper $h_t=1/\sqrt n$ one has $\sum_t h_t^2=n\cdot\tfrac1n=1$, recovering the ordinary periodogram.)

\textbf{(c)} We need $\norm h_2^2=c^2\sum_t g_t^2=1$. With $g=(1,2,1)$,
\[
\sum_{t=1}^{3}g_t^2=1^2+2^2+1^2=6
\quad\Longrightarrow\quad
c=\frac{1}{\sqrt6},\qquad
h=\frac{1}{\sqrt6}\,(1,2,1) .
\]
(Check: $\tfrac16(1+4+1)=1$.) \emph{Tapering vs.\ multitapering.} The two cure different defects of the raw periodogram. \textbf{Tapering} replaces the Fej\'er kernel $\Fej_n$ by the sharper, lower-side-lobe kernel $\abs{H(f)}^2$, reducing \emph{bias} (spectral leakage). It does \emph{nothing} for the variance: a single tapered periodogram still has $\Var\approx\sdf(f)^2$ and is inconsistent. \textbf{Multitapering} averages $K$ tapered periodograms from orthonormal tapers, reducing the \emph{variance} to $\approx\sdf(f)^2/K$ and restoring consistency (let $K\to\infty$ slowly), at the cost of a modest bias increase.

\end{document}
